\subsection{Calibration}\label{\refsec:calibration}
This first model is used to match the case of Argentina only. The full procedure and a description of the data can be found in the main paper. The main application in the paper uses constant parameter values across $t$ for almost all parameters. Thus, when we drop the $t$-subscript in the following, we assume constant parameters.  

\smalltitle{Data.} 
We follow previous literature and let $t$ represent 30-year intervals (this helps with the assumption that capital fully depreciates between time steps). We find population data including projections to determine $\nu_t$ from 1950 to 2070 and amend the data with $t = 1920, 2100$  with constant extrapolation.

We use a household survey to split up the population into $J+1 = 5$ household types. This household survey provides shares $\gamma$, hours worked, and income. In addition, we have data on voting shares represented by $\mu_j$ in the political problem and survival rates $(p_j)$ (where $p_i = p$ is common across $i>0$ types). 

We use the capital income share to identify $\alpha$, set the labour supply elasticity at a baseline value $\xi = 0.3$, take the capital--output ratio to be 3.23 --- Argentina's ratio in the Penn World Table at the baseline year $t=2010$, where the remaining targets are measured too --- and the pension tax rate $\tau_t$ to be 12.5\% in the baseline year. Because a period is thirty years, the ratio the data report is thirty times the model's own $K_t/Y_t$, which is a stock over a single period's output; \cref{\refnumsec:calibration} sets out how that is imposed. 

We use OECD data to define ratio of replacement rates $RR$ at half mean wages and mean wages (represented by groups $j=1, j=3$ in our data). This will help us identify $\theta$ (we do not work with endogenous pension characteristics for Argentina). 

In the Argentina case, we exploit a pension reform around $t=2010$ that drastically changed the degree of universal pensions $\epsilon$. This does not, however, yield a static calibration target, as we see below.


\smalltitle{Calibration.}
Several of the model parameters are set directly using the data. This covers the parts that are not set directly, but have to solved in one way or another.

First, we calibrate the income distribution and relative working hours across the income distribution using $\eta_j$ and $X_j$. Specifically, let $z_j^{\eta}$ and $z_j^{x}$ denote income and working hours of type $j$ relative to the average \emph{formal} household. The reference is the same for all types, informal ones included, and the informal type is compared to it without being part of it, so that
\begin{align}\label{\refeq:calibration:z}
    \sum_{i}\gamma_i z_i^{\eta} = \sum_{i}\gamma_i z_i^{x} = 1,
\end{align}
where, as throughout, the sums over $i$ run over the formal types only. In the data this means dividing each group's income and hours by the population-weighted mean over the formal groups.

With $A_t\equiv w_t(1-\tau_t)+\theta_{t+1}\overline{b}_{t+1}$, formal labour supply \eqref{\refeq:formalOpt} is $h_{t,i} = (\eta_i/X_i)^{\xi}A_t^{\xi}$, so that $h_t = \Gamma_hA_t^{\xi}$ and $h_{t,i}/h_t = (\eta_i/X_i)^{\xi}/\Gamma_h$. The relative income and hours of the formal types are then
\begin{align*}
        z_i^{\eta} &= \dfrac{w_t(1-\tau_t)h_{t,i} \eta_i}{w_t(1-\tau_t)h_t} = \dfrac{\eta_i^{1+\xi}/X_i^{\xi}}{\sum_{j} \gamma_ j\eta_j^{1+\xi}/X_j^{\xi}},  \\
        z_i^{x} &= \frac{h_{t,i}}{\sum_j\gamma_j h_{t,j}} = \frac{(\eta_i/X_i)^{\xi}}{\sum_j \gamma_j(\eta_j/X_j)^{\xi}}.
\end{align*}
Both conditions are ratios and leave two normalizations free. The first is the level of the labour aggregator $\Gamma_h$. The second is the unit in which hours are counted: the rescaling $\eta_i\rightarrow\eta_i/\mu$, $X_i\rightarrow\mu^{-(1+\xi)/\xi}X_i$ leaves $\eta_i^{1+\xi}/X_i^{\xi}$ and therefore efficiency units $\eta_ih_{t,i}$, the aggregate $h_t$ and the equilibrium unchanged, and scales every $h_{t,i}$ by $\mu$. We spend both normalizations in the calibration stage:
\begin{align}\label{\refeq:calibration:yNorm}
    \Gamma_h\equiv\sum_i\gamma_i\frac{\eta_i^{1+\xi}}{X_i^{\xi}} = 1, \qquad \sum_i\gamma_i\left(\frac{\eta_i}{X_i}\right)^{\xi} = 1.
\end{align}
The second one is a real restriction on the informal calibration below, not a convenience. The aggregate $h_t$ weights hours by productivity, and $h_t/\sum_i\gamma_ih_{t,i} = \Gamma_h/\sum_i\gamma_i(\eta_i/X_i)^{\xi}$. Under \eqref{\refeq:calibration:yNorm}, and only under it, $h_t$ equals average formal hours in the same clock unit as the individual $h_{t,i}$ and $h_{1,t}^0$, which lets the informal hours target be stated against $h_t$.

Using $y_i^{\eta} = \eta_i^{1+\xi}/X_i^{\xi}$ and $y_i^x = (\eta_i/X_i)^{\xi}$ as shorthand, we can set up the conditions as linear systems instead:
\begin{align*}
    \left(\bm{z}^{\eta} \bm{\gamma}'-\bm{I}\right)\bm{y}^{\eta} &= \bm{0} \\
    \left(\bm{z}^{x} \bm{\gamma}'-\bm{I}\right) \bm{y}^{x} &= \bm{0},
\end{align*}
where the bold notation indicates vectors and matrices over the formal types. The $\bm{y}$ vectors are eigenvectors with eigenvalue $\bm{\gamma}'\bm{z} = 1$. Since $\bm{z}\bm{\gamma}'$ has rank one, the eigenvector is $\bm{z}$ itself, and \eqref{\refeq:calibration:yNorm} fixes its scale at $\bm{\gamma}'\bm{y} = 1$, so $\bm{y}^{\eta} = \bm{z}^{\eta}$ and $\bm{y}^{x} = \bm{z}^{x}$. The vectors of productivity and disutility of labour follow as
\begin{align}
    \bm{\eta} &= \frac{\bm{y}^{\eta}}{\bm{y}^x} \label{\refeq:calibration:etai} \\
    \bm{X} &= \frac{\bm{\eta}}{(\bm{y}^x)^{1/\xi}}, \label{\refeq:calibration:Xi}
\end{align}
with division and powers taken element by element.

To capture the relative income and working hours of the informal/hand-to-mouth households, requires a bit more. By \eqref{\refeq:calibration:z} and \eqref{\refeq:calibration:yNorm}, the ratios of young informal income and working hours are
\begin{align*}
    z_0^{\eta} &= \frac{w_t^0h_{1,t}^0\eta_{0}}{w_t(1-\tau_t)h_t}, \\
    z_0^x &= \dfrac{h_{1,t}^0}{\sum_i \gamma_{i} h_{t,i}} = \frac{h_{1,t}^0}{h_t}.
\end{align*}
Write $k_t\equiv s_{t-1}/\nu_t$. In equilibrium $h_t = \Theta_{h,t}k_t^{\alpha\xi/(1+\alpha\xi)}$, so the factor prices \eqref{\refeq:factorPrices} and \eqref{\refeq:w0} give $w_t = (1-\alpha)\Theta_{h,t}^{-\alpha}k_t^{\alpha/(1+\alpha\xi)}$ and $w_t^0 = \left((1-\alpha)/\Gamma_{h,t}^{\alpha}\right)^{1/(1+\alpha\xi)}k_t^{\alpha/(1+\alpha\xi)}$. The ratio $z_0^{\eta}/z_0^x = \eta_0w_t^0/\left(w_t(1-\tau_t)\right)$ then identifies $\eta_0$, and the informal labour supply condition \eqref{\refeq:informalOpt}, $h_{1,t}^0 = (\eta_0w_t^0/X_0)^{\xi} = z_0^x\Theta_{h,t}k_t^{\alpha\xi/(1+\alpha\xi)}$, identifies $X_0$. Here $k_t$ drops out of both:
\begin{align}
    \eta_{0} &= \frac{z_0^{\eta}}{z_0^x}\frac{(1-\alpha)(1-\tau_t)}{\Theta_{h,t}^{\alpha}\left(\frac{1-\alpha}{\Gamma_{h,t}^{\alpha}}\right)^{\frac{1}{1+\alpha\xi}}} \label{\refeq:calibration:eta0}\\
    X_{0} &= \eta_{0}\frac{\left(\frac{1-\alpha}{\Gamma_{h,t}^{\alpha}}\right)^{\frac{1}{1+\alpha\xi}}}{\left(\Theta_{h,t}z_0^x\right)^{1/\xi}} \label{\refeq:calibration:X0}
\end{align}
We note that all of this is evaluated at $t=2010$, our baseline year. However, while we can identify most of the terms, the coefficient function $\Theta_{h,t}$ still has to be solved for us to get this. Thus, we add these two conditions determining $\eta_0, X_0$ as targets that we have to solve alongside the politico-economic equilibrium (to get $\Theta_{h,t}$ in the baseline year). 



Given vectors of productivities and disutility of labour for formal households, $\bm{\eta}_i, \bm{X}_i$, we can determine $\theta$ from the ratio of replacement rates. In particular, we solve for $\theta$ (constant across $t$ in this calibration) from the requirement:
 \begin{align*}
        RR \equiv \dfrac{\theta+(1-\theta)\frac{h_t}{\eta_3 h_{t,3}}}{\theta+(1-\theta)\frac{h_t}{\eta_1 h_{t,1}}} = \frac{\theta + (1-\theta)\frac{\Gamma_{h,t} X_3^{\xi}}{\eta_3^{1+\xi}}}{\theta+(1-\theta)\frac{\Gamma_{h,t} X_1^{\xi}}{\eta_1^{1+\xi}}},
\end{align*}
where we have data on $RR$ from OECD data, and the two groups $j=1, 3$ represent half-mean and mean wages in our data for Argentina. Given that we know $\eta_i, X_i$ for formal households already, we can simply compute $\theta$ from
\begin{align}
    \theta = \frac{RR \frac{\Gamma_{h,t} X_1^{\xi}}{\eta_1^{1+\xi}} - \frac{\Gamma_{h,t} X_3^{\xi}}{\eta_3^{1+\xi}}}{1-\frac{\Gamma_{h,t} X_3^{\xi}}{\eta_3^{1+\xi}}-RR\left(1-\frac{\Gamma_{h,t} X_1^{\xi}}{\eta_1^{1+\xi}}\right)}
\end{align}



Finally, we need to determine parameters $\omega$ and discount parameters $\beta$. Note that we assume that they are identical across time and household types in the calibration. They are identified by ensuring that the politico-economic equilibrium replicates a pension tax of 12.5\% and a capital--output ratio of 3.23 in the baseline year 2010. Capital at $t$ is the previous generation's savings, $K_t = s_{t-1}/\nu_t$, so with Cobb--Douglas production the ratio is
\begin{align*}
    \frac{K_t}{Y_t} = \frac{s_{t-1}/\nu_t}{\left(s_{t-1}/\nu_t\right)^{\alpha} h_t^{1-\alpha}} = \frac{\alpha}{R_t},
\end{align*}
a stock over \emph{one period's} output; multiplying by the thirty years in a period puts it in the units the data are reported in. The savings rate
\begin{align*}
    \text{Savings rate}_t = \frac{s_t}{\left(s_{t-1}/\nu_t\right)^{\alpha} h_t^{1-\alpha}}
\end{align*}
is reported alongside it, but is not a target: with capital fully depreciating between periods it equals $K_{t+1}/Y_t$, which is again a stock over a period's flow rather than the annual saving rate the national accounts measure.


Most parameters can be calibrated initially, assuming values for parameters $\xi, \rho$. There are, however, four parameters that cannot be calibrated without knowing the politico-economic equilibrium path: $\omega, \beta, \eta_0, X_0$. As the politico-economic equilibrium cannot be solved without these parameters, the final bit of calibration is solved as a nested fixed point problem. In \cref{\refnumsec} we describe this calibration procedure in more detail. 